\section{问题重述}

药材在热风烘干过程中同时发生升温和脱水，内部状态具有明显的径向差异。题目给出烘房温度、环境水分浓度、药材初始状态、半径变化及多组经验物性关系，要求逐步求解预热阶段分布、变物性热质场、全域含水率达标时间，并分析收缩和物性变化对结果的影响。

四个子问题构成递进关系：问题一建立固定尺寸的基准场；问题二将热物性和扩散率改为状态相关形式；问题三把场变量转化为全域阈值事件；问题四引入给定的收缩半径轨迹，比较几何和整组物性变化的作用。问题一至问题四均以题目给定的初边值和附件时序数据为输入，输出温度或含水率的径向分布以及相应的验证指标。

本文将附件中的烘房温度和环境浓度按时间插值，将长度固定的圆柱药材近似为一维径向问题；问题四进一步用材料坐标将收缩区域映射到固定区间。模型结果只在所声明的有效表面平衡浓度、无显式潜热和均匀收缩等条件下成立。
